% Scope: Spin-half quantum mechanics, thermal ensembles and the classical connection.
% Status: foundational chapter with explicit density-matrix dynamics.
% Prerequisites: Chapter 1 and the front-matter sign conventions.
% Planned figures: Implemented Zeeman levels and Bloch-vector geometry.
% Planned challenge problems: Implemented; solutions are in Appendix E.
\chapter{Quantum Spin Physics}
\label{ch:02}

Nuclear spin is intrinsic angular momentum, not a charged sphere whose
mechanical rotation must be inferred. Its state space, observables and
dynamical law determine what an MRI experiment can prepare and measure.
For uncoupled spin-half nuclei these rules admit an exact three-component
representation of the density operator. This explains why a classical-looking
magnetization vector describes many MRI experiments, even though thermal
polarization and the magnetic moment are quantum mechanical. The operator
framework used here is standard NMR theory \cite{levitt2008}.

\section{Spin as intrinsic angular momentum}
\label{sec:02:spin-as-intrinsic-angular-momentum}
\subsection{States, measurement bases and superposition}
Choose eigenvectors of $\op S_z$ as an orthonormal basis:
$\ket{\uparrow}=(1,0)^{\mathsf T}$ and
$\ket{\downarrow}=(0,1)^{\mathsf T}$.
A pure state is $\ket\psi=a\ket{\uparrow}+b\ket{\downarrow}$,
where $\abs a^2+\abs b^2=1$. A common phase multiplying both coefficients
changes no measurement probability. Their relative phase does:
transverse-basis measurements depend on $a^*b$.

Up to a common phase,
\[
 \ket\psi=\cos(\vartheta/2)\ket{\uparrow}
       +\ee^{\ii\varphi}\sin(\vartheta/2)\ket{\downarrow}.
\]
A $z$ measurement gives $+\hbar/2$ with probability
$\cos^2(\vartheta/2)$ and $-\hbar/2$ with probability
$\sin^2(\vartheta/2)$. Neither outcome implies a pre-existing classical
vector with simultaneous definite Cartesian components. Repeated
measurements on identically prepared systems determine expectation
values, which can be represented by a vector.

\subsection{Pauli matrices and spin operators}
The Pauli matrices and dimensional spin operators are
\begin{equation}
 \sigma_x=\begin{pmatrix}0&1\\1&0\end{pmatrix},\quad
 \sigma_y=\begin{pmatrix}0&-\ii\\\ii&0\end{pmatrix},\quad
 \sigma_z=\begin{pmatrix}1&0\\0&-1\end{pmatrix},\qquad
 \op S_j=\frac{\hbar}{2}\sigma_j .
 \label{eq:02:pauli}
\end{equation}
Direct multiplication gives
$\sigma_j\sigma_k=\delta_{jk}\mat I+
\ii\sum_l\epsilon_{jkl}\sigma_l$ and hence
$\op{\vect S}^{\,2}=3\hbar^2\mat I/4$.
An individual component eigenvalue has magnitude $\hbar/2$;
the squared total angular momentum has eigenvalue
$\hbar^2s(s+1)$ with $s=1/2$. Confusing these magnitudes gives an
incorrect magnetic moment.

For any unit vector $\vect n$,
$(\vect n\cdot\vect\sigma)^2=\mat I$, so the two projectors are
$\mat P_\pm=(\mat I\pm\vect n\cdot\vect\sigma)/2$.
Their expectation values supply measurement probabilities along an
arbitrary axis without explicitly constructing its eigenvectors.

\subsection{Commutators, uncertainty and eigenstates}
The multiplication law implies
$[\op S_j,\op S_k]=\ii\hbar\sum_l\epsilon_{jkl}\op S_l$.
Positivity of the norm of a linear combination of
$(\op A-\langle A\rangle)\ket\psi$ and
$(\op B-\langle B\rangle)\ket\psi$ gives
$\Delta A\,\Delta B\geq\tfrac12\abs{\langle[\op A,\op B]\rangle}$.
In $\ket{\uparrow}$, $\Delta S_x=\Delta S_y=\hbar/2$:
their product saturates the bound $\hbar^2/4$.
Even a fully polarized spin has transverse measurement uncertainty.

The states $\ket{\pm x}=(\ket{\uparrow}\pm\ket{\downarrow})/\sqrt2$
have equal $z$ populations but opposite transverse expectation values.
Similarly, $\ket{+y}=(\ket{\uparrow}+\ii\ket{\downarrow})/\sqrt2$.
Equal longitudinal populations do not imply zero transverse moment.

\section{Magnetic moments and Zeeman splitting}
\label{sec:02:magnetic-moments-and-zeeman-splitting}
\subsection{Gyromagnetic ratio and Hamiltonian}
The magnetic moment operator is $\op{\vect\mu}=\gamma\op{\vect S}$.
In a classical applied field,
\begin{equation}
 \op H(t)=-\gamma\op{\vect S}\cdot\vect B(t)
         =-\frac{\hbar\gamma}{2}\vect B(t)\cdot\vect\sigma .
 \label{eq:02:zeeman}
\end{equation}
The nuclear structure is incorporated in $\gamma$; it is not determined
by modeling the proton as a classical point charge. The units are energy
since $\hbar\gamma B$ is an action times inverse time.

\subsection{Energy levels and transition frequencies}
For positive $\gamma$ and field $B_0\vect e_z$,
\[
 E_\uparrow=-\hbar\omega_0/2,\qquad E_\downarrow=+\hbar\omega_0/2,
 \qquad \Delta E=\hbar\omega_0=hf_0 .
\]
A moment along the field has lower energy. A transverse perturbation
has off-diagonal matrix elements in this basis and drives transitions;
a longitudinal perturbation changes the splitting. Resonance matches
the drive to the splitting, including shielding and local-field shifts.

\begin{figure}[tb]
 \centering\begin{tikzpicture}[x=1cm,y=1cm,font=\small]
\draw[mri axis] (0,0)--(0,3.8) node[above] {$E$};
\draw[mri signal] (.4,.7)--(3,.7);
\draw[mri signal] (.4,3)--(3,3);
\node[anchor=west] at (.5,.35) {$\ket{\uparrow},\ -\hbar\omega_0/2$};
\node[anchor=west] at (.5,3.35) {$\ket{\downarrow},\ +\hbar\omega_0/2$};
\draw[<->,thick] (1.7,.8)--(1.7,2.9) node[midway,right] {RF: $\hbar\omega_0$};
\begin{scope}[xshift=7cm,yshift=1.7cm]
\draw[gray] (0,0) circle (1.55);
\draw[mri guide] (0,0) ellipse (1.55 and .42);
\draw[mri axis] (0,-1.7)--(0,1.9) node[above] {$p_z$};
\draw[mri axis] (-1.8,0)--(1.95,0) node[right] {$p_x$};
\draw[mri axis] (.8,.7)--(-1.3,-1.1) node[below] {$p_y$};
\draw[mri rf,-{Stealth}] (0,0)--(.85,1.15) node[right] {$\vect p$};
\node at (0,-2) {$\norm{\vect p}\leq1$};
\end{scope}
\end{tikzpicture}

 \caption{Positive-$\gamma$ Zeeman levels and a spin-half Bloch vector.
 RF couples the levels. The vector represents expectation values, not
 simultaneous sharp Cartesian spin components.}
 \label{fig:02:zeeman}
\end{figure}

\subsection{Positive and negative gyromagnetic ratios}
For negative $\gamma$, the ordering of the $S_z$ energies and the
precession handedness both reverse. The transition-frequency magnitude
is $\abs\gamma B_0/(2\pi)$. Equilibrium magnetization still points along
the field in this independent-spin paramagnetic model: the preferred
spin orientation and its conversion to magnetic moment both reverse.
A susceptibility proportional to $\gamma^2$ remains positive.
Our proton RF conventions require translation for negative-$\gamma$
nuclei.

\section{Unitary time evolution}
\label{sec:02:unitary-time-evolution}
\subsection{Schrodinger evolution and propagators}
Schrodinger evolution is
$\ii\hbar\partial_t\ket\psi=\op H(t)\ket\psi$.
For constant $\op H$, the solution is
$U(t)\ket{\psi(0)}$ with $U=\exp(-\ii\op Ht/\hbar)$.
Thus
\begin{equation}
 U_0(t)=\begin{pmatrix}
 \ee^{+\ii\omega_0t/2}&0\\0&\ee^{-\ii\omega_0t/2}
 \end{pmatrix}.
 \label{eq:02:free-u}
\end{equation}
The signs follow from the negative magnetic-dipole Hamiltonian.
A $2\pi$ rotation changes a spinor's sign, but returns its density matrix
and single-spin observables to their original values. It does not reverse
the MRI magnetization after a $360^\circ$ RF rotation.

\subsection{Relative phase and Larmor precession}
For amplitudes $a,b$,
$\langle S_x\rangle=\hbar\operatorname{Re}(a^*b)$,
$\langle S_y\rangle=\hbar\operatorname{Im}(a^*b)$ and
$\langle S_z\rangle=\hbar(\abs a^2-\abs b^2)/2$.
Since $a^*(t)b(t)=a^*(0)b(0)\ee^{-\ii\omega_0t}$,
\begin{equation}
 \langle S_x+\ii S_y\rangle(t)
 =\langle S_x+\ii S_y\rangle(0)\ee^{-\ii\omega_0t}.
 \label{eq:02:coherence}
\end{equation}
An initial $\ket{+x}$ state moves toward $-y$.
This is an independent physical check on the precession sign.

\subsection{Expectation values and the classical torque equation}
Use $\dd\langle\op S_i\rangle/\dd t
=(\ii/\hbar)\langle[\op H,\op S_i]\rangle$ and the commutators to obtain
\begin{equation}
 \frac{\dd}{\dd t}\langle\op{\vect S}\rangle
       =\gamma\langle\op{\vect S}\rangle\times\vect B .
 \label{eq:02:ehrenfest}
\end{equation}
Closure is exact because the Hamiltonian is linear in spin.
Multiplication by spin density and $\gamma$ gives the torque part of
the Bloch equations. A coupled-spin term such as
$\op S_{1z}\op S_{2z}$ instead makes first-moment evolution depend on
correlations. A single three-component vector then need not close.

\section{Density operators and ensembles}
\label{sec:02:density-operators-and-ensembles}
\subsection{Pure states, mixtures and coherence}
For a mixture with classical probabilities $p_a$,
$\op\rho=\sum_a p_a\ket{\psi_a}\bra{\psi_a}$.
It is Hermitian, positive semidefinite and trace one.
The mean of $\op A$ is $\Tr(\op\rho\op A)$.
Different ensemble decompositions can give the same density matrix;
single-spin measurements cannot distinguish them.

A pure state has $\op\rho^2=\op\rho$ and $\Tr(\op\rho^2)=1$.
An equal incoherent up/down mixture is $\mat I/2$, whereas
$\ket{+x}\bra{+x}$ has off-diagonal entries $1/2$.
Both have equal $z$ populations. Off-diagonal coherence is basis-dependent;
it is not an extra population of some third spin species.

\subsection{Bloch-vector representation and positivity}
Every Hermitian trace-one $2\times2$ matrix can be written
\begin{equation}
 \op\rho=\frac12(\mat I+\vect p\cdot\vect\sigma)
 =\frac12\begin{pmatrix}1+p_z&p_x-\ii p_y\\p_x+\ii p_y&1-p_z\end{pmatrix}.
 \label{eq:02:density-bloch}
\end{equation}
Here $p_j=\Tr(\op\rho\sigma_j)$ is dimensionless polarization and
$\langle\op{\vect S}\rangle=\hbar\vect p/2$.
The eigenvalues are $(1\pm\norm{\vect p})/2$, using
$(\vect p\cdot\vect\sigma)^2=\norm{\vect p}^2\mat I$.
Physicality requires $\norm{\vect p}\leq1$.
The Bloch sphere surface is pure; its interior is mixed, with purity
$(1+\norm{\vect p}^2)/2$.

For local spin density $\rho_s$,
\begin{equation}
 \vect M=\rho_s\gamma\frac{\hbar}{2}\vect p,\qquad
 M_+=\rho_s\gamma\hbar\,\rho_{21}.
 \label{eq:02:magnetization}
\end{equation}
The \emph{lower} off-diagonal matrix entry encodes $M_+$ in our convention.
These relations apply to a representative local one-spin density matrix,
including an ensemble-averaged reduced state.

\subsection{Dephasing versus loss of single-spin coherence}
Unitary evolution obeys
$\dot{\op\rho}=-(\ii/\hbar)[\op H,\op\rho]$ and preserves eigenvalues
and purity. Frequency dispersion can nevertheless reduce ensemble
transverse mean:
$p_+(t)=p_+(0)\int P(\Delta\omega)\ee^{-\ii\Delta\omega t}\dd\Delta\omega$.
Each member retains magnitude while their vector sum shrinks.
Offset knowledge or a refocusing pulse can undo this cancellation.

Environmental coupling can reduce coherence in a spin's reduced density
matrix too, and this need not be undone by a simple pulse.
Some slowly varying environmental perturbations are refocusable.
The relevant distinctions are correlation time and controllability;
dephasing is not synonymous with an irrecoverable process in every model.

\section{Thermal equilibrium and polarization}
\label{sec:02:thermal-equilibrium-and-polarization}
\subsection{Partition function and Boltzmann populations}
For independent spins at bath temperature $T_{\rm bath}$, let
$\beta=1/(k_{\rm B}T_{\rm bath})$ and $x=\beta\hbar\omega_0/2$.
Then
\[
 Z=\Tr(\ee^{-\beta\op H_0})=2\cosh x,\qquad
 \op\rho_{\rm eq}=\frac{\ee^{x\sigma_z}}{2\cosh x}
                =\frac12(\mat I+\tanh x\,\sigma_z).
\]
The probabilities are $\ee^{\pm x}/(2\cosh x)$, giving
\begin{equation}
 p_\uparrow-p_\downarrow=\tanh x,\qquad
 p_\uparrow/p_\downarrow=\ee^{2x}.
 \label{eq:02:boltzmann}
\end{equation}
The population difference divided by total population differs from the
fractional excess relative to one population. Their small-$x$ limits
differ by a factor of two.

\subsection{High-temperature expansion and its accuracy}
For $\abs x\ll1$, $\tanh x=x-x^3/3+O(x^5)$.
Replacing it by $x$ has relative error approximately $x^2/3$.
Physiological temperature is high relative to the proton Zeeman splitting,
not necessarily to every microscopic energy in tissue.
The identity part of the thermal density matrix greatly exceeds the
polarized deviation. Unitary RF evolution leaves the identity unchanged;
it rotates only that small deviation.

\subsection{Spin density and macroscopic magnetization}
Combining the population and moment expressions gives
\begin{equation}
 M_0=\rho_s\frac{\gamma\hbar}{2}
       \tanh\left(\frac{\hbar\gamma B_0}{2k_{\rm B}T_{\rm bath}}\right)
 \simeq \frac{\rho_s\gamma^2\hbar^2B_0}{4k_{\rm B}T_{\rm bath}} .
 \label{eq:02:curie}
\end{equation}
Units are moment per volume, or A/m. The linear field scaling assumes
independent spins and small polarization. Received voltage also depends
on frequency and coil response; SNR does not follow from $M_0$ alone.

For $N$ independent spins, mean population imbalance is $N\tanh x$
and its standard deviation is $\sqrt{N(1-\tanh^2x)}$.
Macroscopic MRI volumes therefore have large absolute net polarization
despite its small fraction. This is distinct from the question of whether
receiver noise is dominated by spin fluctuations or electrical noise.

\section{Driven spins and the classical limit}
\label{sec:02:driven-spins-and-the-classical-limit}
\subsection{RF interaction Hamiltonian}
For a co-rotating field of amplitude $B_a$ and phase $\phi$, the
rotating-frame Hamiltonian under the rotating-wave approximation is
\[
 \op H_{\rm rot}=-\frac{\hbar}{2}\left[
 \Delta\omega\sigma_z+\omega_1(\cos\phi\,\sigma_x+\sin\phi\,\sigma_y)\right],
 \qquad\omega_1=\gamma B_a .
\]
\Cref{ch:04} derives the transformation and $B_a=\abs{\Btx}$.
On resonance, for constant phase, the propagator is
$U=\cos(\alpha/2)\mat I+\ii\sin(\alpha/2)\vect n\cdot\vect\sigma$,
with $\alpha=\omega_1t$ and $\vect n=(\cos\phi,\sin\phi,0)$.
Separate even and odd exponential powers to derive this formula.

An initially spin-up pure state then has spin-down probability
$\sin^2(\alpha/2)$. A thermal state instead has its polarization rotated:
for $\phi=0$,
$\vect p=(0,p_0\sin\alpha,p_0\cos\alpha)$.
A $90^\circ$ pulse equalizes populations while maximizing transverse
magnetization; population imbalance alone no longer characterizes it.

\subsection{Rotating-wave approximation and resonance}
A linear RF oscillation has equal co- and counter-rotating components.
In a resonant frame one is slow; the other oscillates near twice the carrier.
When drive strength, detuning and envelope variation rates are small
relative to the carrier, the latter averages to a small correction.
A linearly polarized peak field $B_{\rm pk}$ therefore gives
$\omega_1=\gamma B_{\rm pk}/2$. For an ideal purely co-rotating field,
discarding a counter-rotating term is unnecessary. Strong drives or
rapid modulation require the full Hamiltonian.

\subsection{Ensemble Bloch dynamics and hidden microscopic physics}
For a single spin with a linear Hamiltonian the Bloch vector exactly
represents the density matrix, including mixed states.
Relaxation requires an open-system model. In a Markovian spin-half
model with energy-downward and energy-upward rates $W_d,W_u$ and
pure-dephasing rate $\Gamma_\phi$,
\[
 R_1=W_d+W_u,\qquad R_2=R_1/2+\Gamma_\phi .
\]
For protons the lower-energy state is spin up. Detailed balance fixes
the equilibrium polarization. The bound $T_2\leq2T_1$ follows for this
model; empirical effective times in heterogeneous multi-pool data need
not describe this same pair of microscopic processes.

Classical ensemble magnetization suffices when first moments close under
the dynamics. It hides quantum measurement uncertainty, purity and
higher correlations. Scalar couplings in spectroscopy can transfer
information into correlations outside a one-vector model. Routine imaging
often admits local Bloch vectors with relaxation and transport as an
appropriate reduced description.

\section{Worked examples}
\label{sec:02:worked-examples}
\subsection{Derivation and quantitative calculation}
\begin{workedexample}[Thermal polarization]
Use $\gammabar=\SI{42.57748}{\mega\hertz\per\tesla}$,
$\hbar=\SI{1.054571817e-34}{\joule\second}$ and
$k_{\rm B}=\SI{1.380649e-23}{\joule\per\kelvin}$.
At $T_{\rm bath}=\SI{310}{\kelvin}$:
\begin{center}
\begin{tabular}{@{}rrr@{}}\toprule
$B_0$ (T)&$f_0$ (MHz)&$10^6(p_\uparrow-p_\downarrow)$\\\midrule
1.5&63.8662&4.9437\\3.0&127.7324&9.8874\\7.0&298.0424&23.0706\\\bottomrule
\end{tabular}
\end{center}
The rounded free-proton value illustrates scale, not a scanner calibration;
shielding slightly changes actual resonance.
At 7 T the relative high-temperature error is only
$x^2/3\simeq1.77\times10^{-10}$.
\end{workedexample}
\begin{workedexample}[Net moment in a voxel]
For illustrative water-like proton density
$\rho_s=\SI{6.69e28}{\per\cubic\metre}$ at 3 T and 310 K,
$M_0=\SI{9.33e-3}{\ampere\per\metre}$.
A \SI{1}{\cubic\milli\metre} voxel has $6.69\times10^{19}$ protons,
mean population excess $6.62\times10^{14}$ and equilibrium magnetic
moment $\SI{9.33e-12}{\ampere\square\metre}$.
The density is a stated model input, not a universal tissue constant.
\end{workedexample}
\subsection{Assumptions, limiting cases and scanner interpretation}
\begin{workedexample}[Equal populations with coherence]
Apply $U=\ee^{+\ii\pi\sigma_x/4}$ to
$\op\rho_{\rm eq}=(\mat I+p_0\sigma_z)/2$.
Multiplication gives $\op\rho'=(\mat I+p_0\sigma_y)/2$:
both diagonal entries are $1/2$, while $\rho'_{21}=\ii p_0/2$.
Thus $M_+=\ii M_0$. Replacing this state by an incoherent 50:50
mixture would incorrectly predict no signal.
\end{workedexample}

\section{Challenge problems}
\label{sec:02:challenge-problems}
\subsection{Derivation and quantitative design}
\begin{challengeproblem}[A proposed density matrix]
\label{prob:02:tomography}
For $\op\rho=\tfrac12\begin{pmatrix}1+a&b-\ii c\\b+\ii c&1-a\end{pmatrix}$,
derive physicality, purity and Cartesian-axis measurement probabilities.
Evaluate $(a,b,c)=(0.6,0.3,0.4)$ and explain the failure of
$(0.9,0.6,0)$. Distinguish a nonphysical fit from a negative noisy datum.
\end{challengeproblem}
\begin{challengeproblem}[Population conventions]
\label{prob:02:thermal}
Derive the exact spin-half partition function, polarization and magnetization
for signed $\gamma$. Show that equilibrium magnetization is parallel to
$B_0$ for either sign. At 3 T and 310 K compare
$(N_\uparrow-N_\downarrow)/(N_\uparrow+N_\downarrow)$ with
$(N_\uparrow-N_\downarrow)/N_\downarrow$ for protons.
Estimate when linearized polarization has one-percent relative error.
\end{challengeproblem}
\begin{challengeproblem}[A two-pulse interferometer]
\label{prob:02:ramsey}
Starting from $p_0\vect e_z$, apply a phase-zero $\pi/2$ pulse,
free rotating-frame evolution for $\tau$ with detuning $\Delta\omega$,
then a $\pi/2$ pulse of phase $\phi$. Neglect relaxation.
Derive final $p_z$ by quantum propagators and by Bloch rotations.
\end{challengeproblem}
\subsection{Model criticism and integrated reasoning}
\begin{challengeproblem}[Reversible ensemble dephasing]
\label{prob:02:dephase}
Equal subensembles start in $\ket{+x}$ with offsets $+\Omega$ and
$-\Omega$. Find their average density matrix and purity versus time.
Apply a phase-zero $\pi$ pulse at $\tau$ and determine the state at
$2\tau$. Can loss of ensemble purity prove irreversible decoherence?
\end{challengeproblem}
\begin{challengeproblem}[A physical relaxation map]
\label{prob:02:rates}
For energy-downward and energy-upward rates $W_d,W_u$, derive equilibrium
polarization and $T_1$ from the population master equation.
If coherence decays at $(W_d+W_u)/2+\Gamma_\phi$, find $T_2$ and
the allowed ratio $T_2/T_1$. Why does fitting two arbitrary exponentials
not identify a valid microscopic model automatically?
\end{challengeproblem}
